% -----------------------------------------------------------------------------
% Personal customisations: extra packages, structure slides, helper macros.
% Loaded from main.tex with % -----------------------------------------------------------------------------
% Personal customisations: extra packages, structure slides, helper macros.
% Loaded from main.tex with % -----------------------------------------------------------------------------
% Personal customisations: extra packages, structure slides, helper macros.
% Loaded from main.tex with % -----------------------------------------------------------------------------
% Personal customisations: extra packages, structure slides, helper macros.
% Loaded from main.tex with \input{customize.tex}.
%
% Nothing here may change the theme's fonts.  Loading lmodern (or any other
% font package) after the theme silently overrides the Caladea/Carlito pairing
% the theme sets up, which is why it is not loaded.
% -----------------------------------------------------------------------------

\usepackage{amsmath,amssymb}
\usepackage{bm}
\usepackage{booktabs}
\usepackage{listings}

% References are set with \slidecite (free-form, no bibliography, no biber).

% -----------------------------------------------------------------------------
% Structure slides
% -----------------------------------------------------------------------------

% Table of contents: one line per section, subsections hidden.
\setbeamertemplate{section in toc}{$\blacktriangleright$~\inserttocsection}
\setbeamertemplate{section in toc shaded}[default][30]
\setbeamertemplate{subsection in toc}{}

% Frames broken over several slides are not numbered "(cont.)".
\setbeamertemplate{frametitle continuation}{}

% A standalone contents slide, with no section highlighted.  Use it once, after
% the title slide; \AtBeginSection takes over from there.
\newcommand{\outlineslide}{%
  \begingroup
  \setbeamertemplate{footline}{}
  \setbeamertemplate{headline}{}
  \themecolor{mint}
  \begin{frame}[plain]{Outline}
    \tableofcontents
  \end{frame}
  \endgroup
}

% Each section opens with the outline, current section highlighted.
\AtBeginSection[]{%
  \begingroup
  \setbeamertemplate{footline}{}
  \setbeamertemplate{headline}{}
  \themecolor{mint}
  \begin{frame}[plain]{Outline}
    \tableofcontents[currentsection]
  \end{frame}
  \endgroup
}

% Each subsection opens with its own divider slide.
\AtBeginSubsection[]{%
  \begingroup
  \setbeamertemplate{footline}{}
  \setbeamertemplate{headline}{}
  \begin{frame}[plain,c]
    \centering
    {\usebeamerfont{framesubtitle}\usebeamercolor[fg]{framesubtitle}%
     \smallcapsline{Section \thesection.\thesubsection}\par}
    \vspace{1.5ex}
    {\usebeamerfont{title}\usebeamercolor[fg]{title}\subsecname\par}
  \end{frame}
  \endgroup
}

% -----------------------------------------------------------------------------
% Code listings: Julia, terminal sessions and prompts to an AI agent
% -----------------------------------------------------------------------------
% Use them in frames marked [fragile]:
%   \begin{julia} ... \end{julia}        Julia source
%   \begin{terminal} ... \end{terminal}  shell commands and their output
%   \begin{prompt} ... \end{prompt}      what you type to the agent
%   \begin{agentlog} ... \end{agentlog}  what an agent does (schematic transcript)
% Greek letters and a few symbols in the source are mapped to maths, so Julia
% code can be pasted verbatim from DrivenSpin/ (pdfLaTeX has no Unicode mono font).
\definecolor{codegray}{rgb}{0.5,0.5,0.5}
\definecolor{codepurple}{rgb}{0.58,0,0.82}

\lstdefinelanguage{Julia}{%
  morekeywords={abstract,begin,break,catch,const,continue,do,else,elseif,end,
    export,false,finally,for,function,global,if,import,in,let,local,macro,
    module,mutable,quote,return,struct,true,try,using,while,where},
  morecomment=[l]\#,
  morecomment=[n]{\#=}{=\#},
  morestring=[b]",
  sensitive=true,
}

\lstdefinestyle{unibasbase}{
    basicstyle=\ttfamily\scriptsize,
    breakatwhitespace=false,
    breaklines=true,
    keepspaces=true,
    columns=fullflexible,
    showspaces=false,
    showstringspaces=false,
    showtabs=false,
    tabsize=4,
    upquote=true,
    extendedchars=true,
    literate=
      {Ω}{{$\Omega$}}1 {Δ}{{$\Delta$}}1 {Γ}{{$\Gamma$}}1 {σ}{{$\sigma$}}1
      {ρ}{{$\rho$}}1 {ψ}{{$\psi$}}1 {π}{{$\pi$}}1 {φ}{{$\phi$}}1
      {γ}{{$\gamma$}}1 {μ}{{$\mu$}}1 {ω}{{$\omega$}}1 {ħ}{{$\hbar$}}1
      {≈}{{$\approx$}}1 {±}{{$\pm$}}1 {≤}{{$\leq$}}1 {≥}{{$\geq$}}1 {→}{{$\to$}}1
      {₁}{{$_1$}}1 {₂}{{$_2$}}1 {ₑ}{{$_e$}}1 {†}{{$^\dagger$}}1
      {⟩}{{$\rangle$}}1 {⟨}{{$\langle$}}1 {└}{{$\llcorner$}}1 {●}{{\textbullet}}1
      {·}{{$\cdot$}}1 {—}{{---}}1 {✓}{{$\checkmark$}}1 {✗}{{$\times$}}1 {×}{{$\times$}}1,
    xleftmargin=4pt,
    xrightmargin=4pt,
    aboveskip=4pt,
    belowskip=4pt,
    frame=single,
    framerule=0.4pt,
    framesep=4pt,
}

\lstdefinestyle{unibasel}{
    style=unibasbase,
    language=Julia,
    commentstyle=\color{unibasPetrol},
    keywordstyle=\color{unibasRed},
    stringstyle=\color{codepurple},
    rulecolor=\color{unibasLightGrey},
    backgroundcolor=\color{unibasPaper},
}
\lstset{style=unibasel}

\lstnewenvironment{julia}[1][]{\lstset{style=unibasel,#1}}{}

\lstnewenvironment{terminal}[1][]{\lstset{style=unibasbase,language={},
    rulecolor=\color{unibasAnthracite},
    backgroundcolor=\color{unibasAnthracite!6},
    commentstyle=\color{unibasGrey},
    morecomment=[l]\#,
    #1}}{}

\lstnewenvironment{prompt}[1][]{\lstset{style=unibasbase,language={},
    basicstyle=\ttfamily\footnotesize,
    rulecolor=\color{unibasPetrol},
    backgroundcolor=\color{unibasMint!20},
    title={\hfill\tiny\color{unibasPetrol}\smallcapsline{prompt}},
    #1}}{}

\lstnewenvironment{agentlog}[1][]{\lstset{style=unibasbase,language={},
    rulecolor=\color{unibasGrey},
    backgroundcolor=\color{unibasWhite},
    title={\hfill\tiny\color{unibasGrey}\smallcapsline{agent transcript (schematic)}},
    #1}}{}

% -----------------------------------------------------------------------------
% Helper macros
% -----------------------------------------------------------------------------

% Coloured hyperlinks
\newcommand{\hrefcol}[2]{\textcolor{unibasPetrol}{\href{#1}{#2}}}
\newcommand{\hlinkcol}[1]{\hrefcol{#1}{#1}}

% Centred, narrower paragraph (a statement that should not run edge to edge)
\newcommand{\centerstate}[1]{%
  \centering
  \begin{columns}
    \begin{column}{0.8\textwidth}
      #1
    \end{column}
  \end{columns}%
}

% Coloured emphasis.  \c... is Basel red (the highlight colour),
% \b... is petrol (structure).  Use sparingly: on a slide, one emphasis wins.
\newcommand{\ctextbf}[1]{\textbf{\textcolor{unibasRed}{#1}}}
\newcommand{\btextbf}[1]{\textbf{\textcolor{unibasPetrol}{#1}}}
\newcommand{\ctextsl}[1]{\textsl{\textcolor{unibasRed}{#1}}}
\newcommand{\btextsl}[1]{\textsl{\textcolor{unibasPetrol}{#1}}}
\newcommand{\cemph}[1]{\emph{\textcolor{unibasRed}{#1}}}
\newcommand{\bemph}[1]{\emph{\textcolor{unibasPetrol}{#1}}}

% About / acknowledgement page
\newcommand{\aboutpage}[2]{%
  \begingroup
  \setbeamertemplate{footline}{}
  \setbeamertemplate{headline}{}
  \themecolor{mint}
  \begin{frame}[plain,c]{#1}
    \centering
    \begin{minipage}{0.85\textwidth}
      \usebeamercolor[fg]{normal text}
      \centering #2
    \end{minipage}
  \end{frame}
  \endgroup
}
.
%
% Nothing here may change the theme's fonts.  Loading lmodern (or any other
% font package) after the theme silently overrides the Caladea/Carlito pairing
% the theme sets up, which is why it is not loaded.
% -----------------------------------------------------------------------------

\usepackage{amsmath,amssymb}
\usepackage{bm}
\usepackage{booktabs}
\usepackage{listings}

% References are set with \slidecite (free-form, no bibliography, no biber).

% -----------------------------------------------------------------------------
% Structure slides
% -----------------------------------------------------------------------------

% Table of contents: one line per section, subsections hidden.
\setbeamertemplate{section in toc}{$\blacktriangleright$~\inserttocsection}
\setbeamertemplate{section in toc shaded}[default][30]
\setbeamertemplate{subsection in toc}{}

% Frames broken over several slides are not numbered "(cont.)".
\setbeamertemplate{frametitle continuation}{}

% A standalone contents slide, with no section highlighted.  Use it once, after
% the title slide; \AtBeginSection takes over from there.
\newcommand{\outlineslide}{%
  \begingroup
  \setbeamertemplate{footline}{}
  \setbeamertemplate{headline}{}
  \themecolor{mint}
  \begin{frame}[plain]{Outline}
    \tableofcontents
  \end{frame}
  \endgroup
}

% Each section opens with the outline, current section highlighted.
\AtBeginSection[]{%
  \begingroup
  \setbeamertemplate{footline}{}
  \setbeamertemplate{headline}{}
  \themecolor{mint}
  \begin{frame}[plain]{Outline}
    \tableofcontents[currentsection]
  \end{frame}
  \endgroup
}

% Each subsection opens with its own divider slide.
\AtBeginSubsection[]{%
  \begingroup
  \setbeamertemplate{footline}{}
  \setbeamertemplate{headline}{}
  \begin{frame}[plain,c]
    \centering
    {\usebeamerfont{framesubtitle}\usebeamercolor[fg]{framesubtitle}%
     \smallcapsline{Section \thesection.\thesubsection}\par}
    \vspace{1.5ex}
    {\usebeamerfont{title}\usebeamercolor[fg]{title}\subsecname\par}
  \end{frame}
  \endgroup
}

% -----------------------------------------------------------------------------
% Code listings: Julia, terminal sessions and prompts to an AI agent
% -----------------------------------------------------------------------------
% Use them in frames marked [fragile]:
%   \begin{julia} ... \end{julia}        Julia source
%   \begin{terminal} ... \end{terminal}  shell commands and their output
%   \begin{prompt} ... \end{prompt}      what you type to the agent
%   \begin{agentlog} ... \end{agentlog}  what an agent does (schematic transcript)
% Greek letters and a few symbols in the source are mapped to maths, so Julia
% code can be pasted verbatim from DrivenSpin/ (pdfLaTeX has no Unicode mono font).
\definecolor{codegray}{rgb}{0.5,0.5,0.5}
\definecolor{codepurple}{rgb}{0.58,0,0.82}

\lstdefinelanguage{Julia}{%
  morekeywords={abstract,begin,break,catch,const,continue,do,else,elseif,end,
    export,false,finally,for,function,global,if,import,in,let,local,macro,
    module,mutable,quote,return,struct,true,try,using,while,where},
  morecomment=[l]\#,
  morecomment=[n]{\#=}{=\#},
  morestring=[b]",
  sensitive=true,
}

\lstdefinestyle{unibasbase}{
    basicstyle=\ttfamily\scriptsize,
    breakatwhitespace=false,
    breaklines=true,
    keepspaces=true,
    columns=fullflexible,
    showspaces=false,
    showstringspaces=false,
    showtabs=false,
    tabsize=4,
    upquote=true,
    extendedchars=true,
    literate=
      {Ω}{{$\Omega$}}1 {Δ}{{$\Delta$}}1 {Γ}{{$\Gamma$}}1 {σ}{{$\sigma$}}1
      {ρ}{{$\rho$}}1 {ψ}{{$\psi$}}1 {π}{{$\pi$}}1 {φ}{{$\phi$}}1
      {γ}{{$\gamma$}}1 {μ}{{$\mu$}}1 {ω}{{$\omega$}}1 {ħ}{{$\hbar$}}1
      {≈}{{$\approx$}}1 {±}{{$\pm$}}1 {≤}{{$\leq$}}1 {≥}{{$\geq$}}1 {→}{{$\to$}}1
      {₁}{{$_1$}}1 {₂}{{$_2$}}1 {ₑ}{{$_e$}}1 {†}{{$^\dagger$}}1
      {⟩}{{$\rangle$}}1 {⟨}{{$\langle$}}1 {└}{{$\llcorner$}}1 {●}{{\textbullet}}1
      {·}{{$\cdot$}}1 {—}{{---}}1 {✓}{{$\checkmark$}}1 {✗}{{$\times$}}1 {×}{{$\times$}}1,
    xleftmargin=4pt,
    xrightmargin=4pt,
    aboveskip=4pt,
    belowskip=4pt,
    frame=single,
    framerule=0.4pt,
    framesep=4pt,
}

\lstdefinestyle{unibasel}{
    style=unibasbase,
    language=Julia,
    commentstyle=\color{unibasPetrol},
    keywordstyle=\color{unibasRed},
    stringstyle=\color{codepurple},
    rulecolor=\color{unibasLightGrey},
    backgroundcolor=\color{unibasPaper},
}
\lstset{style=unibasel}

\lstnewenvironment{julia}[1][]{\lstset{style=unibasel,#1}}{}

\lstnewenvironment{terminal}[1][]{\lstset{style=unibasbase,language={},
    rulecolor=\color{unibasAnthracite},
    backgroundcolor=\color{unibasAnthracite!6},
    commentstyle=\color{unibasGrey},
    morecomment=[l]\#,
    #1}}{}

\lstnewenvironment{prompt}[1][]{\lstset{style=unibasbase,language={},
    basicstyle=\ttfamily\footnotesize,
    rulecolor=\color{unibasPetrol},
    backgroundcolor=\color{unibasMint!20},
    title={\hfill\tiny\color{unibasPetrol}\smallcapsline{prompt}},
    #1}}{}

\lstnewenvironment{agentlog}[1][]{\lstset{style=unibasbase,language={},
    rulecolor=\color{unibasGrey},
    backgroundcolor=\color{unibasWhite},
    title={\hfill\tiny\color{unibasGrey}\smallcapsline{agent transcript (schematic)}},
    #1}}{}

% -----------------------------------------------------------------------------
% Helper macros
% -----------------------------------------------------------------------------

% Coloured hyperlinks
\newcommand{\hrefcol}[2]{\textcolor{unibasPetrol}{\href{#1}{#2}}}
\newcommand{\hlinkcol}[1]{\hrefcol{#1}{#1}}

% Centred, narrower paragraph (a statement that should not run edge to edge)
\newcommand{\centerstate}[1]{%
  \centering
  \begin{columns}
    \begin{column}{0.8\textwidth}
      #1
    \end{column}
  \end{columns}%
}

% Coloured emphasis.  \c... is Basel red (the highlight colour),
% \b... is petrol (structure).  Use sparingly: on a slide, one emphasis wins.
\newcommand{\ctextbf}[1]{\textbf{\textcolor{unibasRed}{#1}}}
\newcommand{\btextbf}[1]{\textbf{\textcolor{unibasPetrol}{#1}}}
\newcommand{\ctextsl}[1]{\textsl{\textcolor{unibasRed}{#1}}}
\newcommand{\btextsl}[1]{\textsl{\textcolor{unibasPetrol}{#1}}}
\newcommand{\cemph}[1]{\emph{\textcolor{unibasRed}{#1}}}
\newcommand{\bemph}[1]{\emph{\textcolor{unibasPetrol}{#1}}}

% About / acknowledgement page
\newcommand{\aboutpage}[2]{%
  \begingroup
  \setbeamertemplate{footline}{}
  \setbeamertemplate{headline}{}
  \themecolor{mint}
  \begin{frame}[plain,c]{#1}
    \centering
    \begin{minipage}{0.85\textwidth}
      \usebeamercolor[fg]{normal text}
      \centering #2
    \end{minipage}
  \end{frame}
  \endgroup
}
.
%
% Nothing here may change the theme's fonts.  Loading lmodern (or any other
% font package) after the theme silently overrides the Caladea/Carlito pairing
% the theme sets up, which is why it is not loaded.
% -----------------------------------------------------------------------------

\usepackage{amsmath,amssymb}
\usepackage{bm}
\usepackage{booktabs}
\usepackage{listings}

% References are set with \slidecite (free-form, no bibliography, no biber).

% -----------------------------------------------------------------------------
% Structure slides
% -----------------------------------------------------------------------------

% Table of contents: one line per section, subsections hidden.
\setbeamertemplate{section in toc}{$\blacktriangleright$~\inserttocsection}
\setbeamertemplate{section in toc shaded}[default][30]
\setbeamertemplate{subsection in toc}{}

% Frames broken over several slides are not numbered "(cont.)".
\setbeamertemplate{frametitle continuation}{}

% A standalone contents slide, with no section highlighted.  Use it once, after
% the title slide; \AtBeginSection takes over from there.
\newcommand{\outlineslide}{%
  \begingroup
  \setbeamertemplate{footline}{}
  \setbeamertemplate{headline}{}
  \themecolor{mint}
  \begin{frame}[plain]{Outline}
    \tableofcontents
  \end{frame}
  \endgroup
}

% Each section opens with the outline, current section highlighted.
\AtBeginSection[]{%
  \begingroup
  \setbeamertemplate{footline}{}
  \setbeamertemplate{headline}{}
  \themecolor{mint}
  \begin{frame}[plain]{Outline}
    \tableofcontents[currentsection]
  \end{frame}
  \endgroup
}

% Each subsection opens with its own divider slide.
\AtBeginSubsection[]{%
  \begingroup
  \setbeamertemplate{footline}{}
  \setbeamertemplate{headline}{}
  \begin{frame}[plain,c]
    \centering
    {\usebeamerfont{framesubtitle}\usebeamercolor[fg]{framesubtitle}%
     \smallcapsline{Section \thesection.\thesubsection}\par}
    \vspace{1.5ex}
    {\usebeamerfont{title}\usebeamercolor[fg]{title}\subsecname\par}
  \end{frame}
  \endgroup
}

% -----------------------------------------------------------------------------
% Code listings: Julia, terminal sessions and prompts to an AI agent
% -----------------------------------------------------------------------------
% Use them in frames marked [fragile]:
%   \begin{julia} ... \end{julia}        Julia source
%   \begin{terminal} ... \end{terminal}  shell commands and their output
%   \begin{prompt} ... \end{prompt}      what you type to the agent
%   \begin{agentlog} ... \end{agentlog}  what an agent does (schematic transcript)
% Greek letters and a few symbols in the source are mapped to maths, so Julia
% code can be pasted verbatim from DrivenSpin/ (pdfLaTeX has no Unicode mono font).
\definecolor{codegray}{rgb}{0.5,0.5,0.5}
\definecolor{codepurple}{rgb}{0.58,0,0.82}

\lstdefinelanguage{Julia}{%
  morekeywords={abstract,begin,break,catch,const,continue,do,else,elseif,end,
    export,false,finally,for,function,global,if,import,in,let,local,macro,
    module,mutable,quote,return,struct,true,try,using,while,where},
  morecomment=[l]\#,
  morecomment=[n]{\#=}{=\#},
  morestring=[b]",
  sensitive=true,
}

\lstdefinestyle{unibasbase}{
    basicstyle=\ttfamily\scriptsize,
    breakatwhitespace=false,
    breaklines=true,
    keepspaces=true,
    columns=fullflexible,
    showspaces=false,
    showstringspaces=false,
    showtabs=false,
    tabsize=4,
    upquote=true,
    extendedchars=true,
    literate=
      {Ω}{{$\Omega$}}1 {Δ}{{$\Delta$}}1 {Γ}{{$\Gamma$}}1 {σ}{{$\sigma$}}1
      {ρ}{{$\rho$}}1 {ψ}{{$\psi$}}1 {π}{{$\pi$}}1 {φ}{{$\phi$}}1
      {γ}{{$\gamma$}}1 {μ}{{$\mu$}}1 {ω}{{$\omega$}}1 {ħ}{{$\hbar$}}1
      {≈}{{$\approx$}}1 {±}{{$\pm$}}1 {≤}{{$\leq$}}1 {≥}{{$\geq$}}1 {→}{{$\to$}}1
      {₁}{{$_1$}}1 {₂}{{$_2$}}1 {ₑ}{{$_e$}}1 {†}{{$^\dagger$}}1
      {⟩}{{$\rangle$}}1 {⟨}{{$\langle$}}1 {└}{{$\llcorner$}}1 {●}{{\textbullet}}1
      {·}{{$\cdot$}}1 {—}{{---}}1 {✓}{{$\checkmark$}}1 {✗}{{$\times$}}1 {×}{{$\times$}}1,
    xleftmargin=4pt,
    xrightmargin=4pt,
    aboveskip=4pt,
    belowskip=4pt,
    frame=single,
    framerule=0.4pt,
    framesep=4pt,
}

\lstdefinestyle{unibasel}{
    style=unibasbase,
    language=Julia,
    commentstyle=\color{unibasPetrol},
    keywordstyle=\color{unibasRed},
    stringstyle=\color{codepurple},
    rulecolor=\color{unibasLightGrey},
    backgroundcolor=\color{unibasPaper},
}
\lstset{style=unibasel}

\lstnewenvironment{julia}[1][]{\lstset{style=unibasel,#1}}{}

\lstnewenvironment{terminal}[1][]{\lstset{style=unibasbase,language={},
    rulecolor=\color{unibasAnthracite},
    backgroundcolor=\color{unibasAnthracite!6},
    commentstyle=\color{unibasGrey},
    morecomment=[l]\#,
    #1}}{}

\lstnewenvironment{prompt}[1][]{\lstset{style=unibasbase,language={},
    basicstyle=\ttfamily\footnotesize,
    rulecolor=\color{unibasPetrol},
    backgroundcolor=\color{unibasMint!20},
    title={\hfill\tiny\color{unibasPetrol}\smallcapsline{prompt}},
    #1}}{}

\lstnewenvironment{agentlog}[1][]{\lstset{style=unibasbase,language={},
    rulecolor=\color{unibasGrey},
    backgroundcolor=\color{unibasWhite},
    title={\hfill\tiny\color{unibasGrey}\smallcapsline{agent transcript (schematic)}},
    #1}}{}

% -----------------------------------------------------------------------------
% Helper macros
% -----------------------------------------------------------------------------

% Coloured hyperlinks
\newcommand{\hrefcol}[2]{\textcolor{unibasPetrol}{\href{#1}{#2}}}
\newcommand{\hlinkcol}[1]{\hrefcol{#1}{#1}}

% Centred, narrower paragraph (a statement that should not run edge to edge)
\newcommand{\centerstate}[1]{%
  \centering
  \begin{columns}
    \begin{column}{0.8\textwidth}
      #1
    \end{column}
  \end{columns}%
}

% Coloured emphasis.  \c... is Basel red (the highlight colour),
% \b... is petrol (structure).  Use sparingly: on a slide, one emphasis wins.
\newcommand{\ctextbf}[1]{\textbf{\textcolor{unibasRed}{#1}}}
\newcommand{\btextbf}[1]{\textbf{\textcolor{unibasPetrol}{#1}}}
\newcommand{\ctextsl}[1]{\textsl{\textcolor{unibasRed}{#1}}}
\newcommand{\btextsl}[1]{\textsl{\textcolor{unibasPetrol}{#1}}}
\newcommand{\cemph}[1]{\emph{\textcolor{unibasRed}{#1}}}
\newcommand{\bemph}[1]{\emph{\textcolor{unibasPetrol}{#1}}}

% About / acknowledgement page
\newcommand{\aboutpage}[2]{%
  \begingroup
  \setbeamertemplate{footline}{}
  \setbeamertemplate{headline}{}
  \themecolor{mint}
  \begin{frame}[plain,c]{#1}
    \centering
    \begin{minipage}{0.85\textwidth}
      \usebeamercolor[fg]{normal text}
      \centering #2
    \end{minipage}
  \end{frame}
  \endgroup
}
.
%
% Nothing here may change the theme's fonts.  Loading lmodern (or any other
% font package) after the theme silently overrides the Caladea/Carlito pairing
% the theme sets up, which is why it is not loaded.
% -----------------------------------------------------------------------------

\usepackage{amsmath,amssymb}
\usepackage{bm}
\usepackage{booktabs}
\usepackage{listings}

% References are set with \slidecite (free-form, no bibliography, no biber).

% -----------------------------------------------------------------------------
% Structure slides
% -----------------------------------------------------------------------------

% Table of contents: one line per section, subsections hidden.
\setbeamertemplate{section in toc}{$\blacktriangleright$~\inserttocsection}
\setbeamertemplate{section in toc shaded}[default][30]
\setbeamertemplate{subsection in toc}{}

% Frames broken over several slides are not numbered "(cont.)".
\setbeamertemplate{frametitle continuation}{}

% A standalone contents slide, with no section highlighted.  Use it once, after
% the title slide; \AtBeginSection takes over from there.
\newcommand{\outlineslide}{%
  \begingroup
  \setbeamertemplate{footline}{}
  \setbeamertemplate{headline}{}
  \themecolor{mint}
  \begin{frame}[plain]{Outline}
    \tableofcontents
  \end{frame}
  \endgroup
}

% Each section opens with the outline, current section highlighted.
\AtBeginSection[]{%
  \begingroup
  \setbeamertemplate{footline}{}
  \setbeamertemplate{headline}{}
  \themecolor{mint}
  \begin{frame}[plain]{Outline}
    \tableofcontents[currentsection]
  \end{frame}
  \endgroup
}

% Each subsection opens with its own divider slide.
\AtBeginSubsection[]{%
  \begingroup
  \setbeamertemplate{footline}{}
  \setbeamertemplate{headline}{}
  \begin{frame}[plain,c]
    \centering
    {\usebeamerfont{framesubtitle}\usebeamercolor[fg]{framesubtitle}%
     \smallcapsline{Section \thesection.\thesubsection}\par}
    \vspace{1.5ex}
    {\usebeamerfont{title}\usebeamercolor[fg]{title}\subsecname\par}
  \end{frame}
  \endgroup
}

% -----------------------------------------------------------------------------
% Code listings: Julia, terminal sessions and prompts to an AI agent
% -----------------------------------------------------------------------------
% Use them in frames marked [fragile]:
%   \begin{julia} ... \end{julia}        Julia source
%   \begin{terminal} ... \end{terminal}  shell commands and their output
%   \begin{prompt} ... \end{prompt}      what you type to the agent
%   \begin{agentlog} ... \end{agentlog}  what an agent does (schematic transcript)
% Greek letters and a few symbols in the source are mapped to maths, so Julia
% code can be pasted verbatim from DrivenSpin/ (pdfLaTeX has no Unicode mono font).
\definecolor{codegray}{rgb}{0.5,0.5,0.5}
\definecolor{codepurple}{rgb}{0.58,0,0.82}

\lstdefinelanguage{Julia}{%
  morekeywords={abstract,begin,break,catch,const,continue,do,else,elseif,end,
    export,false,finally,for,function,global,if,import,in,let,local,macro,
    module,mutable,quote,return,struct,true,try,using,while,where},
  morecomment=[l]\#,
  morecomment=[n]{\#=}{=\#},
  morestring=[b]",
  sensitive=true,
}

\lstdefinestyle{unibasbase}{
    basicstyle=\ttfamily\scriptsize,
    breakatwhitespace=false,
    breaklines=true,
    keepspaces=true,
    columns=fullflexible,
    showspaces=false,
    showstringspaces=false,
    showtabs=false,
    tabsize=4,
    upquote=true,
    extendedchars=true,
    literate=
      {Ω}{{$\Omega$}}1 {Δ}{{$\Delta$}}1 {Γ}{{$\Gamma$}}1 {σ}{{$\sigma$}}1
      {ρ}{{$\rho$}}1 {ψ}{{$\psi$}}1 {π}{{$\pi$}}1 {φ}{{$\phi$}}1
      {γ}{{$\gamma$}}1 {μ}{{$\mu$}}1 {ω}{{$\omega$}}1 {ħ}{{$\hbar$}}1
      {≈}{{$\approx$}}1 {±}{{$\pm$}}1 {≤}{{$\leq$}}1 {≥}{{$\geq$}}1 {→}{{$\to$}}1
      {₁}{{$_1$}}1 {₂}{{$_2$}}1 {ₑ}{{$_e$}}1 {†}{{$^\dagger$}}1
      {⟩}{{$\rangle$}}1 {⟨}{{$\langle$}}1 {└}{{$\llcorner$}}1 {●}{{\textbullet}}1
      {·}{{$\cdot$}}1 {—}{{---}}1 {✓}{{$\checkmark$}}1 {✗}{{$\times$}}1 {×}{{$\times$}}1,
    xleftmargin=4pt,
    xrightmargin=4pt,
    aboveskip=4pt,
    belowskip=4pt,
    frame=single,
    framerule=0.4pt,
    framesep=4pt,
}

\lstdefinestyle{unibasel}{
    style=unibasbase,
    language=Julia,
    commentstyle=\color{unibasPetrol},
    keywordstyle=\color{unibasRed},
    stringstyle=\color{codepurple},
    rulecolor=\color{unibasLightGrey},
    backgroundcolor=\color{unibasPaper},
}
\lstset{style=unibasel}

\lstnewenvironment{julia}[1][]{\lstset{style=unibasel,#1}}{}

\lstnewenvironment{terminal}[1][]{\lstset{style=unibasbase,language={},
    rulecolor=\color{unibasAnthracite},
    backgroundcolor=\color{unibasAnthracite!6},
    commentstyle=\color{unibasGrey},
    morecomment=[l]\#,
    #1}}{}

\lstnewenvironment{prompt}[1][]{\lstset{style=unibasbase,language={},
    basicstyle=\ttfamily\footnotesize,
    rulecolor=\color{unibasPetrol},
    backgroundcolor=\color{unibasMint!20},
    title={\hfill\tiny\color{unibasPetrol}\smallcapsline{prompt}},
    #1}}{}

\lstnewenvironment{agentlog}[1][]{\lstset{style=unibasbase,language={},
    rulecolor=\color{unibasGrey},
    backgroundcolor=\color{unibasWhite},
    title={\hfill\tiny\color{unibasGrey}\smallcapsline{agent transcript (schematic)}},
    #1}}{}

% -----------------------------------------------------------------------------
% Helper macros
% -----------------------------------------------------------------------------

% Coloured hyperlinks
\newcommand{\hrefcol}[2]{\textcolor{unibasPetrol}{\href{#1}{#2}}}
\newcommand{\hlinkcol}[1]{\hrefcol{#1}{#1}}

% Centred, narrower paragraph (a statement that should not run edge to edge)
\newcommand{\centerstate}[1]{%
  \centering
  \begin{columns}
    \begin{column}{0.8\textwidth}
      #1
    \end{column}
  \end{columns}%
}

% Coloured emphasis.  \c... is Basel red (the highlight colour),
% \b... is petrol (structure).  Use sparingly: on a slide, one emphasis wins.
\newcommand{\ctextbf}[1]{\textbf{\textcolor{unibasRed}{#1}}}
\newcommand{\btextbf}[1]{\textbf{\textcolor{unibasPetrol}{#1}}}
\newcommand{\ctextsl}[1]{\textsl{\textcolor{unibasRed}{#1}}}
\newcommand{\btextsl}[1]{\textsl{\textcolor{unibasPetrol}{#1}}}
\newcommand{\cemph}[1]{\emph{\textcolor{unibasRed}{#1}}}
\newcommand{\bemph}[1]{\emph{\textcolor{unibasPetrol}{#1}}}

% About / acknowledgement page
\newcommand{\aboutpage}[2]{%
  \begingroup
  \setbeamertemplate{footline}{}
  \setbeamertemplate{headline}{}
  \themecolor{mint}
  \begin{frame}[plain,c]{#1}
    \centering
    \begin{minipage}{0.85\textwidth}
      \usebeamercolor[fg]{normal text}
      \centering #2
    \end{minipage}
  \end{frame}
  \endgroup
}
